\section{Conclusion}
\label{sec:conclusion}

We treated clothing-fit prediction as what it structurally is -- an ordinal
classification problem over three ordered levels -- rather than defaulting
to an unordered multinomial classifier or collapsing it into a binary
signal, and paired the fitted model with split conformal prediction so its
output is a calibrated set, not a bare label. Both pieces are classical
machine learning: a cumulative link model fit by maximizing an explicit
likelihood with \texttt{scipy.optimize}, and a closed-form calibration
procedure over a held-out partition, implemented from scratch in NumPy. No
neural network, embedding layer, or gradient-descent training loop appears
anywhere in the pipeline.

The ordinal structure measurably helped: against an unordered multinomial
logistic regression trained on identical features, the proportional-odds
model nearly halved the mean absolute ordinal error ($0.3535$ versus
$0.6298$) and produced visibly tighter conformal sets ($1.72$ versus $2.08$
classes on average) at the same target coverage. Against the naive
majority-class floor, it traded a small amount of macro F1 generosity to
the majority class for a model that actually uses the shopper's and
garment's features, reflected in more than double the majority baseline's
macro F1.

\subsection{Limitations}

\textbf{Synthetic fixture.} All numbers in this report are computed on the
synthetic fixture, not the real Kaggle dataset; \texttt{scripts/download\_data.py
--kaggle} provides the path to real data when credentials are available, and
the feature pipeline, model, and tests are schema-driven rather than tuned to
the fixture's specific distribution.

\textbf{Proportional-odds assumption.} The shared-weight cumulative link
model assumes every feature's effect on the cumulative log-odds is constant
across thresholds. If, for example, body type affected the \emph{small}
threshold very differently from the \emph{large} threshold, a
non-proportional (threshold-varying) extension would fit better; we did not
find evidence requiring that extra complexity on this fixture, but a
partial-proportional-odds variant is the natural next step on real data.

\textbf{Conformal set size on a 3-class problem.} With only three classes,
an "informative" conformal set can only take sizes $\{1, 2, 3\}$; the method
scales more informatively to problems with more ordinal levels (e.g. a 5-star
or 10-point sizing scale), where a tight conformal set communicates
substantially more than it can here.

\subsection{What a reader should take away}

The point this project makes is not that an ordinal model is marginally more
accurate. It is that treating an explicitly ordered outcome as ordered, and
reporting a model's own uncertainty honestly instead of suppressing it behind
a single confident label, are both small implementation choices with an
outsized effect on whether a recommendation can be trusted.
